\documentclass[10pt,a4paper]{amsart}
\usepackage[margin=19mm]{geometry}
\usepackage{amsmath,amssymb,mathtools}
\usepackage[scr=rsfs,cal=euler]{mathalfa}
\usepackage{xfrac}
\usepackage[unicode,hidelinks,bookmarks=false]{hyperref}
\newcommand{\ie}{i.e.\ }
\newcommand{\cf}{cf.\ }
\newcommand{\resp}{respectively\ }
\newcommand{\Subsection}{\subsection}
\providecommand{\nobreakdash}{\nobreak\hbox{-}}
\setlength{\emergencystretch}{2em}
\multlinegap0pt
\usepackage{amssymb}
\usepackage{float}
\usepackage{enumerate}
\usepackage{bm}
\usepackage{xcolor}
\usepackage[normalem]{ulem}
\newtheorem{prop}{Proposition}
\DeclareMathOperator{\rank}{\text{rank}}
\title{On $U$-unitary Cayley graphs \\ over finite rings}
\author{Tung T. Nguyen, Nguyen Duy T\^an}
\date{Source: arXiv:2603.21239v1 (CC BY 4.0)}
\begin{document}
\maketitle

\noindent\emph{Source and licence.} On $U$-unitary Cayley graphs \\ over finite rings. Version arXiv:2603.21239v1 (CC BY 4.0), distributed by arXiv under the Creative Commons Attribution 4.0 International licence, http://creativecommons.org/licenses/by/4.0/, as recorded in the arXiv licence metadata for this exact version and shown on the abstract page https://arxiv.org/abs/2603.21239v1. Full licence notice: \texttt{licenses/arxiv-2603.21239-CC-BY-4.0.txt}.\par\smallskip

\noindent\textit{Editorial scope.} Counted unit: the article's prop prop:similar-sln (source0.tex lines 622--624). The article's complete original proof of it is delivered with it (source0.tex lines 625--638, one \texttt{proof} environment of 14 lines). No mathematics is written, repaired or re-derived: the statement and the proof are the article's own lines, in the article's own notation, and no retained line is substituted or re-flowed.  No further source-local context is needed: every cross-reference of every retained line resolves inside the two delivered spans.  Nothing was omitted from any retained span, so the package carries no omission record.  No retained line contains a citation, so the wrapper carries no bibliography and no bibliography entry.  No retained line contains a figure environment, an image-inclusion command or drawing code, so no image is delivered and the package declares no asset dependency.  Editorial formatting only, in addition: an AMS \texttt{amsart} preamble that restates the article's own declarations and macros used by the retained lines, namely the environments \texttt{prop} on separate counters, together with the standard packages those lines need.  The retained environments print in this document as Proposition 1 in order of appearance, whereas the article prints the same items with the numbering of its own full document.  The complete native source of the article as distributed in the arXiv e-print for this exact version is bundled unchanged as \texttt{sources/196928/source0.tex}.
\begin{prop} \label{prop:similar-sln}
    $|SL_n(F) \backslash M_n(F) \slash SL_n(F)|=n+|F|-1.$
\end{prop}

\begin{proof}
    For $A,B \in M_n(F)$ such that $A \sim_{SL_n(F)} B$, then $\det(A)=\det(B).$ Therefore, if $A,B$ are invertible and $\det(A) \neq \det(B)$, then $A$ and $B$ are not equivalent. This shows that 
    \[ |SL_n(F) \backslash GL_n(F) \slash SL_n(F)|=|F|^{\times}=|F|-1.\]
    
    
    
    On the other hand, we claim that if $A \sim_{GL_n(F)} B$ and $A,B$ are not invertible, then $A \sim_{SL_n(F)} B.$ Let $r=\rank(A)=\rank(B).$ Then $r<n$. Let $I_r$ be the identity matrix of size $r\times r$ and $\widehat{I_r}$ be the following $n\times n$-matrix 
    \[ 
\widehat{I_r} = I_r\oplus 0 = \begin{bmatrix} I_r & 0 \\ 0 & 0 \end{bmatrix}. \]
By row and column operations, we know that there exists $P,Q \in GL_n(F)$ such that $A=P \widehat{I_r} Q.$ Let $P'$ (respectively $Q'$) be the new matrix obtained by multiplying the last column of $P$ (respectively the last row of $Q$) by $\dfrac{1}{\det(P)}$ (respectively $\dfrac{1}{\det(Q)}$). Then, we still have 
\[ A = P\widehat{I_r} Q = P' \widehat{I_r} Q'.\]
Furthermore, $\det(P')=\det(Q')=1.$ Therefore, $A \sim_{SL_n(F)} \widehat{I_r}.$ Similarly, $B \sim_{SL_n(F)} \widehat{I_r}$, and hence $A \sim_{SL_n(F)} B.$ This shows that the rank map gives an isomorphism between $SL_n(F) \backslash [M_n(F)-GL_n(F)] \slash SL_n(F)$ and the set $\{0, 1, \ldots, n\}.$ In summary, we have  $|SL_n(F) \backslash M_n(F) \slash SL_n(F)|=n+|F|-1.$

\end{proof}

\end{document}
